%!TEX program = lualatex
\documentclass[ngerman,headheight=70pt]{scrartcl}
\usepackage[ngerman]{babel}
\usepackage{blindtext}
\usepackage[utf8]{luainputenc}
\usepackage{scrlayer-scrpage}
\usepackage[T1]{fontenc}
\usepackage{fontspec}
\usepackage{multicol}
\usepackage{enumitem}
\usepackage[german=quotes]{csquotes}
\usepackage[bottom=2cm,footskip=8mm]{geometry}
\setmainfont{Source Sans Pro}
\MakeOuterQuote{"}

\setlength{\parindent}{0pt}
\setlength{\parskip}{6pt}

\pagestyle{plain.scrheadings}
\thispagestyle{scrheadings}
\lohead{\vspace*{-1.5cm}\includegraphics[scale=0.4]{../images/stupa_logo_rot.pdf}}
\rohead{\\Präsidium des\\Studierendenparlaments}

\lofoot[Universität Hamburg · Präsidium des Studierendenparlaments\\\hrule Von-Melle-Park 5 · D-20146 Hamburg]{Universität Hamburg · Präsidium des Studierendenparlaments\\\hrule Von-Melle-Park 5 · D-20146 Hamburg}
\cofoot[]{}
\cohead[\thepage]{}

\renewcommand{\thesubsection}{TOP \arabic{subsection}}
\renewcommand{\thesubsubsection}{\arabic{subsubsection}.}
\setcounter{subsection}{-1}

\begin{document}
    UHH · StuPa-Präsidium · Von-Melle-Park 5 · D-20146 Hamburg

    %%% PAGE ONE (EDIT LATER) %%%%
    % Hier den Titel des Protokolls eintragen
    \section*{Protokoll der 1. Sitzung des Studierendenparlaments vom 14. April 2016}

    % Protokollanten eintragen
    \textbf{Protokoll: Geoffrey N. Youett, Jim Martens}\\
    \textbf{Ort: HWP-Hörsaal}\\
    \textbf{Beginn: 18.25 Uhr}\\
    \textbf{Ende: 0.50 Uhr}

    % Wohl die komplizierteste Stelle im Protokoll
    % Als Faustregel sollten nicht mehr als 3 Namen auf einer Zeile stehen
    % Einfach an dieser Vorlage orientieren
    % Im Idealfall ist hier nicht viel zu ändern, da kaum jemand fehlt.
    % Rücktritte müssen natürlich pro Protokoll angepasst werden
    \vspace{0.5cm}
    \begin{tabular}{ll}
        Anwesend: & \\
            RCDS (5 Sitze): & Ramon Weilinger, Antonia Niecke, Ramin Shakiba, \\
                            & Jennifer Maack, Benjamin Welling \\
             CampusGrün (14 Sitze): & Laura Franzen, Geoffrey Youett, Elena Rysikova, \\
                                   & Philipp Droll, Yasemin Günther, Melf Johannsen,\\
                                   & Tahnee Herzig, Mario Moldenhauer, Jim Martens,\\
                                   & Svenja Horn, Mirzo Ulugbek Khatamov, Armin Günther,\\
                                   & Martin Sievert \\
             Bier-Liste (2 Sitze) : & Jakob Pape \\
             WiWi (2 Sitze): & Claas-Friso Hente, Kay Zöllmer \\
             Unicorns (5 Sitze): & Katharina Kucza, Johannes Peplow, Annkathrin Löffler, \\
                                 & Andreas Hartkamp, Marielle Hermstrüwer \\
             Liste LINKS (3 Sitze): & Gunhild Berdal, Till Petersen, Sinah Mielich \\
             HWP (2 Sitze): & Samet Gunay, Ajdina Karahasan \\
             MIN (4 Sitze): & Ailina Salten, Lotte Rullkötter, Thea Wahlers, Jan Detampel \\
             SDS* (3 Sitze): & Mena Winkler, Jacob Petersein, Artur Brückmann \\
             Bart-LISTE (2 Sitze): & Timo Zeimet, Dominic Laumer \\
             LHG (1 Sitz): & Tobias Heisig \\
             harte zeiten (1 Sitz): & Tobias Berking \\
             Jura (1 Sitz): & \\
             AL (2 Sitze): & Karima Schulze, Henri Weber \\
            & \\
        Entschuldigt: & Freya Schmitz (CampusGrün)\\
                                &\\
        Unentschuldigt abwesend: & Karen Martirosian (Bier), Johann Baumhoefener (Jura) \\
                                &\\
        Rücktritte: & Lasse Kleinlützum (MIN-Liste) \(\rightarrow\) Thea Wahlers \\
                    & Kerstin Riecke (CampusGrün) \(\rightarrow\) Blerta Vila \(\rightarrow\) Martin Sievert\\
                    & Elvis Milojevic (HWP-Liste) \(\rightarrow\) Samet Gunay
    \end{tabular}
    \newpage
    %%% Vorgeschlagene Tagesordnung %%%
    \underline{Vorgeschlagene Tagesordnung}
    \begin{enumerate}[label={\textbf{Top \theenumi}},leftmargin=*]
        % folgende Zeile bei Sitzungsfortsetzungen einkommentieren und x durch
        % die Zahl des zuletzt behandelten TOPs ersetzen
        %\setcounter{enumi}{x}
        \item Geschäftsordnung (V1617-005, V1617-009)
        \item Wahl des StuPa-Präsidiums
        \item RIS-Wahl
            % Beispiel für einen TOP mit 2 Unterpunkten
            \begin{enumerate}
                \item Bestätigung der Wahlniederschrift (V1617-003)
                \item Bestätigung der Referentinnen (V1617-003A01)
            \end{enumerate}
        \item Queer-Wahl
            \begin{enumerate}
                \item Bestätigung der Wahlniederschrift (V1617-006)
                \item Bestätigung der Referent*innen
            \end{enumerate}
        \item Wahl des Satzungs-, Wahlordnungs-, und Geschäftsordnungsausschusses
        \item Wahl des Ausschusses gegen Rechts
        \item Wahl des Haushaltsausschusses
        \item Wahl des Wirtschaftsrats
        \item
            \begin{enumerate}
                \item Verfahren zur Wahl zum Ältestenrat
                \item Wahl des Ältestenrats
            \end{enumerate}
        \item
            \begin{enumerate}
                \item Rechenschaftsbericht des amtierenden AStA
                \item Fragen und Diskussion
                \item Entlastung des AStA
            \end{enumerate}
        \item Wahl des neuen AStA-Vorstandes
            \begin{enumerate}
                \item Diskussion VS-Thesen
                \item Wahl des AStA-Vorstandes
            \end{enumerate}
        \item Bestätigung der AStA-Referent*innen
        \item \#uhh hilft (V1617-007)
        \item Gegen Rechts (V1617-008)
        \item Verschiedenes
    \end{enumerate}

    \newpage

    % Hier beginnt Teil A. Bei Fortsetzungen kommt hier nur der Teil über
    % die Feststellung der endgültigen Fassung der Tagesordnung.
    % Für sonstige Sitzungen ist dies bereits das Gerüst für die Formalia.
    % Man muss nur noch die protokollierten Inhalte eintragen.

    {\Large \underline{Teil A}}

    \subsection{Formalia}

    \subsubsection{Geschäftsbericht Präsidium}

    \subsubsection{Anfragen an das Präsidium}

    \subsubsection{Geschäftsbericht AStA}

    % Beispiel-GO-Handlung
    \textit{Die Überprüfung der Beschlussfähigkeit ergibt, dass das Parlament
    mit 42 anwesenden Parlamentarier*innen beschlussfähig ist.}

    \subsubsection{Anfragen an den AStA}

    \subsubsection{Dringlichkeitsanträge des AStA}

    \subsubsection{Aktuelle Stunde (falls entsprechender Antrag vorliegt)}

    % Bei Anführungsstrichen einfach die normalen auf der Tastatur verwenden.
    % Beim Kompilieren werden diese durch typographische ersetzt.
    Der Antrag "Gegen Rechts" wird im Rahmen der aktuellen Stunde diskutiert.

    \subsubsection{Feststellung der endgültigen Fassung des Teils B der Tagesordnung}

    %%% BESCHLOSSENE TAGESORDNUNG %%%
    \underline{Endgültige Fassung der Tagesordnung}
    \begin{enumerate}[label={\textbf{Top \theenumi}},leftmargin=*]
        % folgende Zeile bei Sitzungsfortsetzungen einkommentieren und x durch
        % die Zahl des zuletzt behandelten TOPs ersetzen
        %\setcounter{enumi}{x}
        \item Gegen Rechts (V1617-008)
        \item Geschäftsordnung (V1617-005, V1617-009)
        \item Wahl des StuPa-Präsidiums
        \item RIS-Wahl
            \begin{enumerate}
                \item Bestätigung der Wahlniederschrift (V1617-003)
                \item Bestätigung der Referentinnen (V1617-003A01)
            \end{enumerate}
        \item Queer-Wahl
            \begin{enumerate}
                \item Bestätigung der Wahlniederschrift (V1617-006)
                \item Bestätigung der Referent*innen
            \end{enumerate}
        \item Wahl des Satzungs-, Wahlordnungs-, und Geschäftsordnungsausschusses
        \item Wahl des Ausschusses gegen Rechts
        \item Wahl des Haushaltsausschusses
        \item Wahl des Wirtschaftsrats
        \item
            \begin{enumerate}
                \item Verfahren zur Wahl zum Ältestenrat
                \item Wahl des Ältestenrats
            \end{enumerate}
        \item \#uhh hilft (V1617-007)
        \item
            \begin{enumerate}
                \item Rechenschaftsbericht des amtierenden AStA
                \item Fragen und Diskussion
                \item Entlastung des AStA
            \end{enumerate}
        \item Wahl des neuen AStA-Vorstandes
            \begin{enumerate}
                \item Diskussion VS-Thesen
                \item Wahl des AStA-Vorstandes
            \end{enumerate}
        \item Bestätigung der AStA-Referent*innen
        \item Verschiedenes
    \end{enumerate}

    \subsubsection{Feststellung der Beschlussfähigkeit}

    \subsubsection{Genehmigung der Protokolle der vorangegangenen Sitzungen}

    \vspace{1cm}
    {\Large \underline{Teil B}}

    % Teil B beginnt
    % Dieser wird sich für jede Sitzung unterscheiden.
    % Jeder TOP ist eine subsection.
    % Vorhandene Beispiele anfangs entfernen und durch tatsächliche TOPs ersetzen

    % Folgende Zeile bei Sitzungsfortsetzungen einkommentieren und x mit
    % der Zahl des zuletzt behandelten TOPs ersetzen.
    %\setcounter{subsection}{x}

    \subsection{Gegen Rechts}

    % Beispiel für einen Verweis auf etwas noch nicht protokolliertes
    % TODO: Anhang anfügen
    Ramon Weilinger stellt einen Ersatzungsantrag zur AfD vor (siehe Anhang).

    \subsection{Wahl des StuPa-Präsidiums}

    Geoffrey Youett und Fabian Schnack übernehmen die Wahlleitung.

    Es werden 2 Listen eingereicht:

    % Beispiel für eine nummerierte Liste bei Wahlen
    \underline{Liste 1}
    \begin{enumerate}
        \item Jim Martens
        \item Gunhild Berdal
        \item Henri Weber
    \end{enumerate}

    \underline{Liste 2}
    \begin{enumerate}
        \item Ramon Weilinger
        \item Ajdina Karahasan
    \end{enumerate}

    Die Kandidierenden stellen sich vor und beantworten Fragen.
    Es werden 41 Stimmzettel ausgegeben.
    Auf Liste 1 entfallen 28 Stimmen, auf Liste 2 entfallen 13 Stimmen.

    Damit sind Jim Martens und Gunhild Berdal von Liste 1 und Ramon Weilinger
    von Liste 2 in das Präsidium gewählt.

    Um sich zu besprechen, schlägt das neue Präsidium eine Pause von 15 Minuten
    bis 23.55 Uhr vor.

    \subsection{Wahl des Satzungs-, Wahlordnungs- und Geschäftsordnungsausschusses}

    % Beispiel für eine nicht nummerierte Liste von Kandidierenden bei Wahlen
    \underline{Liste der Kandidierenden}
    \begin{itemize}
        \item Laura Franzen
        \item Matthias Krupse
        \item Ramon Weilinger
        \item Vincent Orth
        \item Melf Johannsen
        \item Kevin Högy
        \item Artur Brückmann
        \item Merzo Khatamov
        \item Kolja Kolb
    \end{itemize}

    Philipp Droll schlägt vor, dass der SWOGA 9 Mitglieder hat. Es gibt keinen
    Widerspruch.

    Die Kandidierenden gehen nach vorne und stellen sich vor. Es gibt keine
    Fragen.

    Sie werden einstimmig gewählt.

\end{document}
